\uuid{94FF}
\titre{Vecteurs}
\auteur{bodin ; hanani ; mzari}
\organisation{exo7}
\licence{CC BY-NC-SA 4.0}
\datecreate{2020-01-20}
\niveau{L1}
\module{Géométrie}
\chapitre{Géométrie du plan}
\sousChapitre{Distance, norme, produit scalaire}
\difficulte{1}
\theme{vecteurs}
\exo7qcmid{lille-1.geometrie-plane.006}

\contenu{
\begin{question}
\qclassification{140.01,140.02}
\qtags{motcle=vecteurs}
On considère les points $A(0,1)$, $B(2,3)$ et $C(1,1)$. Quelles sont les bonnes réponses ?
\begin{answers}
    \bad{Les droites $(AB)$ et $(OC)$ sont confondues.}
    \bad{Les droites $(AB)$ et $(OC)$ sont perpendiculaires.}
    \good{Les droites $(AB)$ et $(OC)$ sont parallèles.}
    \bad{Les droites $(AB)$ et $(OC)$ sont sécantes.}
\end{answers}
\begin{explanations}
On a $\overrightarrow{AB}=(2,2)=2\overrightarrow{OC}$. Les droites $(AB)$ et $(OC)$ sont parallèles.
\end{explanations}
\end{question}
}
